% report_naver.tex -- Full bilingual EN/KO academic report, Naver Place branch.
% Compile with xelatex (twice, for cross-references).
%
% Format follows 이민영·정동영·김지영 (2026), 소비자학연구 37(1):
%   국문초록/Abstract -> I.서론 -> II.이론적 배경 -> III.연구방법 -> IV.연구결과 -> V.결론 및 시사점 -> 참고문헌
%
% Tables 1-4 (\tblVariables \tblDescriptives \tblCorrelations \tblVIF) are
% generated by 10_01_descriptive_tables.py via tables_naver.tex -- re-run
% that script and recompile this file to refresh them.
% Tables 5-9 (model results) are written by hand from
% logs/pipeline/20261008_1223/09_01_naver_models.log (N=14,401, run 2026-10-08)
% and must be updated by hand if the models are re-run.
%
% Citations are limited to sources actually read in this project:
%   Lee, Jeong & Kim (2026)   -- primary reference, read in full
%   Wang, Karmakar, Wang & Pei (2025) -- read (abstract + introduction)
%   Karabila et al. (2024)    -- read (abstract)
%   Chang (2021); Ko et al. (2018); Bae et al. (2017) -- read (abstract only,
%     used only for the one background sentence their abstracts support)
%   Fresneda & Gefen (2019); Namvar & Chua (2023); Dellarocas (2003) -- NOT
%     read directly; cited only as reported inside Wang et al. (2025) /
%     Lee et al. (2026), marked "as cited in ..." throughout.
% Every other claim that would normally carry a citation is marked
% [SOURCE NEEDED: ...] with a concrete search suggestion instead of a
% fabricated reference -- see the \todo markers below and the note at the
% end of both language halves.
%
% Scope: Naver Place only.

\documentclass[10pt,a4paper]{article}

\usepackage{fontspec}
\usepackage{xeCJK}
\setCJKmainfont{Malgun Gothic}
\setCJKsansfont{Malgun Gothic}
\xeCJKsetup{CJKspace=true}

\usepackage[top=2.0cm, bottom=2.2cm, left=2.2cm, right=2.2cm]{geometry}
\usepackage{booktabs}
\usepackage{amsmath}
\usepackage{enumitem}
\usepackage{microtype}
\usepackage{caption}
\usepackage{float}
\usepackage{longtable}
\usepackage[hidelinks]{hyperref}

\captionsetup[table]{name=표, labelsep=period, font=small, labelfont=bf,
                     skip=4pt, justification=centering}

\linespread{1.04}
\setlength{\parindent}{0pt}
\setlength{\parskip}{5pt}

\newcommand{\sect}[1]{\vspace{10pt}\noindent{\Large\bfseries #1}\par\vspace{3pt}}
\newcommand{\subsect}[1]{\vspace{6pt}\noindent{\large\bfseries #1}\par\vspace{2pt}}
\newcommand{\todo}[1]{\par\noindent\fbox{\parbox{0.96\linewidth}{\small\textit{[SOURCE NEEDED: #1]}}}\par}

% Generated by 10_01_descriptive_tables.py -- do not edit by hand.
% Requires \usepackage{float} in the including document ([H] placement).
% Defines: \tblPlaces \tblVariables \tblDescriptives \tblCorrelations \tblVIF

\newcommand{\tblPlaces}{%
\begin{table}[H]\centering%
\caption{관광지별 리뷰 수 / Reviews per attraction}%
\footnotesize\begin{tabular}{llrrr}\hline%
관광지 & English & Reviews & w/ text & w/ rating \\ \hline%
서울스카이 & Seoul Sky & 1,092 & 1,025 & 792 \\%
남산골한옥마을 & Namsangol Village & 1,075 & 1,033 & 826 \\%
에버랜드 & Everland & 995 & 935 & 851 \\%
N서울타워 & N Seoul Tower & 939 & 817 & 198 \\%
동대문디자인플라자 & DDP & 668 & 453 & 232 \\%
덕수궁 & Deoksugung & 646 & 539 & 72 \\%
운현궁 & Unhyeongung & 616 & 562 & 384 \\%
서대문형무소역사관 & Seodaemun Prison History Hall & 610 & 331 & 27 \\%
창덕궁 & Changdeokgung & 610 & 399 & 38 \\%
올림픽공원 & Olympic Park & 609 & 499 & 104 \\%
해운대해수욕장 & Haeundae Beach & 608 & 457 & 53 \\%
서대문자연사박물관 & Seodaemun Museum & 608 & 320 & 20 \\%
광안리해수욕장 & Gwangalli Beach & 607 & 505 & 77 \\%
국립중앙박물관 & Natl. Museum & 606 & 543 & 221 \\%
여의도한강공원 & Yeouido Park & 605 & 443 & 87 \\%
코엑스 & COEX & 590 & 359 & 198 \\%
창경궁 & Changgyeonggung & 590 & 418 & 26 \\%
경복궁 & Gyeongbokgung & 580 & 477 & 91 \\%
낙산공원 & Naksan Park & 435 & 249 & 108 \\%
북촌한옥마을 & Bukchon Village & 296 & 165 & 118 \\%
북한산 & Bukhansan & 253 & 155 & 26 \\%
롯데월드 & Lotte World & 180 & 85 & 10 \\%
경희궁 & Gyeonghuigung & 170 & 98 & 12 \\%
은평한옥마을 & Eunpyeong Village & 156 & 125 & 36 \\%
인사동길 & Insadong-gil & 109 & 68 & 6 \\%
광장시장 & Gwangjang Mkt & 59 & 32 & 14 \\%
명동성당 & Myeongdong Cathedral & 39 & 33 & 18 \\%
홍대거리 & Hongdae St. & 30 & 22 & 6 \\%
서촌 & Seochon & 20 & 18 & 1 \\%
\hline\end{tabular}%
\end{table}%
}

\newcommand{\tblVariables}{%
\begin{table}[H]\centering%
\caption{변수요약 / Variable summary}%
\footnotesize\begin{tabular}{llp{7.6cm}}\hline%
변수 종류 & 변수 이름 & 변수 설명 \\ \hline%
종속변수 (DV) & Revisit (재방문) & 해당 관광지를 2회차 이상 방문한 작성자의 리뷰 여부 (0/1) \\%
 & Keyword Count (키워드 수) & 리뷰당 작성자가 선택한 키워드 수 (0$\sim$5), 정보 다양성 지표 \\%
 & Rating (별점) & 작성자가 부여한 별점 (1$\sim$5) \\%
 & Reactions (반응 수) & 해당 리뷰에 반응을 남긴 이용자 수, 유용성 지표 \\%
독립변수 (IV) & Companion (동반자) & 방문 동반자 유형 더미: 혼자·연인·친구·아이·가족·동료 \\%
 & Purpose (방문 목적) & 방문 목적 더미: 여행·나들이·일상·데이트·친목 \\%
 & Waited (대기) & 방문 시 대기 경험 여부 (0/1) \\%
 & Reserved (예약) & 사전 예약 여부 (0/1) \\%
 & Sentiment (감성 점수) & 리뷰 텍스트의 감성 점수 (-1 부정 $\sim$ 1 긍정) \\%
 & Tag themes (키워드 주제) & 키워드를 10개 주제로 분류한 더미: 시각·볼거리·혼잡·한적·시설·접근성·가족·가격·음식·물놀이 \\%
통제변수 (Control) & Text Length (텍스트 길이) & 리뷰 텍스트 길이의 로그값 \\%
 & Image Count (사진 수) & 리뷰에 첨부된 사진 수 \\%
 & User Reviews (작성자 리뷰 수) & 작성자가 누적 작성한 리뷰 수의 로그값 \\%
 & User Followers (팔로워 수) & 작성자의 팔로워 수의 로그값 \\%
 & Review Age (리뷰 경과일) & 리뷰 작성 후 경과일의 로그값, 반응 수 모형에서 노출 기간을 통제 \\%
 & Verification (인증) & 영수증 인증 리뷰 여부 (0/1) \\%
 & Season (계절) & 방문 계절 더미: 여름·가을·겨울 (봄 기준) \\%
 & Weekend (주말) & 주말 방문 여부 (0/1) \\%
 & Place FE (관광지 고정효과) & 관광지별 고정효과, 장소 고유 특성 통제 \\%
\hline\end{tabular}\end{table}%
}

\newcommand{\tblDescriptives}{%
\begin{table}[H]\centering%
\caption{기술통계표 / Descriptive statistics}%
\footnotesize\begin{tabular}{lrrrrrr}\hline%
Statistic & N & Mean & St.\ Dev. & Min & Median & Max \\ \hline%
Revisit (재방문) & 14,401 & 0.218 & 0.413 & 0 & 0 & 1 \\%
Keyword Count (키워드 수) & 14,401 & 2.97 & 1.883 & 0 & 3 & 5 \\%
Rating (별점) & 4,652 & 4.682 & 0.732 & 0.5 & 5 & 5 \\%
Reactions (반응 수) & 14,401 & 1.437 & 7.415 & 0 & 0 & 122 \\%
Sentiment (감성 점수) & 11,165 & 0.653 & 0.457 & -0.973 & 0.87 & 0.988 \\%
Text Length (텍스트 길이) & 11,165 & 112.089 & 106.805 & 14 & 71 & 753 \\%
Image Count (사진 수) & 14,401 & 0.235 & 1.412 & 0 & 0 & 20 \\%
User Reviews (작성자 리뷰 수) & 14,401 & 823.26 & 1,009.4 & 1 & 482 & 19,082.0 \\%
Review Age (리뷰 경과일) & 14,401 & 588.324 & 558.142 & 0 & 428 & 3,957.0 \\%
\hline\end{tabular}\end{table}%
}

\newcommand{\tblCorrelations}{%
\begin{table}[H]\centering%
\caption{상관분석표 / Correlation matrix}%
\footnotesize\begin{tabular}{lrrrrrrrrr}\hline%
 & 1 & 2 & 3 & 4 & 5 & 6 & 7 & 8 & 9 \\ \hline%
1. Revisit (재방문) & 1.000 &  &  &  &  &  &  &  &  \\%
2. Keyword Count (키워드 수) & 0.140 & 1.000 &  &  &  &  &  &  &  \\%
3. Rating (별점) & 0.069 & 0.264 & 1.000 &  &  &  &  &  &  \\%
4. Reactions (반응 수) & 0.194 & 0.097 & 0.051 & 1.000 &  &  &  &  &  \\%
5. Sentiment (감성 점수) & -0.067 & 0.116 & 0.514 & -0.039 & 1.000 &  &  &  &  \\%
6. Text Length (텍스트 길이) & 0.198 & 0.204 & 0.000 & 0.221 & -0.095 & 1.000 &  &  &  \\%
7. Image Count (사진 수) & 0.036 & 0.075 & 0.044 & 0.057 & 0.016 & 0.101 & 1.000 &  &  \\%
8. User Reviews (작성자 리뷰 수) & 0.255 & 0.031 & 0.036 & 0.212 & -0.082 & 0.159 & 0.046 & 1.000 &  \\%
9. Review Age (리뷰 경과일) & -0.140 & -0.395 & -0.094 & -0.083 & -0.027 & -0.137 & -0.057 & 0.034 & 1.000 \\%
\hline\end{tabular}\end{table}%
}

\newcommand{\tblVIF}{%
\begin{table}[H]\centering%
\caption{분산팽창계수 / Variance inflation factors}%
\footnotesize\begin{tabular}{lrlrlr}\hline%
변수 / Variable & VIF & 변수 / Variable & VIF & 변수 / Variable & VIF \\ \hline%
has\_visit\_info & 2.82 & comp\_family & 1.21 & sentiment & 1.06 \\%
waited & 1.10 & comp\_colleague & 1.25 & n\_tags & 1.54 \\%
reserved & 1.84 & purp\_travel & 1.38 & revisit & 1.16 \\%
comp\_alone & 1.69 & purp\_outing & 1.65 & ln\_length & 1.10 \\%
comp\_partner & 1.75 & purp\_daily & 1.45 & ln\_user\_reviews & 1.54 \\%
comp\_friend & 1.53 & purp\_date & 1.67 &  &  \\%
comp\_child & 1.70 & purp\_social & 1.39 &  &  \\%
\hline\end{tabular}%
\par\vspace{2pt}\footnotesize\textit{All values below 5; multicollinearity is not a concern. 모든 값이 5 미만으로 다중공선성 문제는 없다.}%
\end{table}%
}

% Shared body for the variable-map table; used by \tblVariableMapEN and \tblVariableMapKO.
% Notation key: $Y_{ij}$ = outcome of that model; $\beta_k$ = coefficient in x_ij vector;
%   $\delta_j$ = place fixed effect; $\gamma$ = ln_review_age exposure offset (M4 only);
%   $\beta_k{\times}N$ = N binary dummies each with their own \beta_k; -- = not included.
% Scale abbreviations used below:
%   bin. = binary (0/1)   ord. = ordinal   cont. = continuous   count = count (integer)   categ. = categorical
\newcommand{\tblVariableMapEN}{%
\begin{table}[H]\centering%
\caption{Variable-to-model mapping}%
\footnotesize\setlength{\tabcolsep}{3.5pt}%
\begin{tabular}{lp{2.7cm}p{2.8cm}lp{1.4cm}p{1.4cm}p{1.4cm}p{1.4cm}}\hline%
Type & Concept (Table~2) & Dataset column(s) & Scale & M1 & M2 & M3 & M4 \\\hline%
\multicolumn{8}{l}{\quad\textit{Dependent variables}}\\\hline%
DV & Revisit             & \texttt{revisit}                          & bin.          & $Y_{ij}$              & $\beta_k$             & $\beta_k$             & $\beta_k$             \\%
DV & Keyword Count       & \texttt{n\_tags}                          & ord.\ 0--5    & $\beta_k$             & $Y_{ij}$              & $\beta_k$             & $\beta_k$             \\%
DV & Rating              & \texttt{rating}                           & ord.\ 1--5    & --                    & --                    & $Y_{ij}$              & --                    \\%
DV & Reactions           & \texttt{reactions}                        & count         & --                    & --                    & --                    & $Y_{ij}$              \\\hline%
\multicolumn{8}{l}{\quad\textit{Independent variables}}\\\hline%
IV & Companion           & \texttt{comp\_*} ($\times$6)              & bin.          & $\beta_k{\times}6$    & $\beta_k{\times}6$    & $\beta_k{\times}4^a$  & $\beta_k{\times}6$    \\%
IV & Purpose             & \texttt{purp\_*} ($\times$5)              & bin.          & $\beta_k{\times}5$    & $\beta_k{\times}5$    & $\beta_k{\times}2^a$  & $\beta_k{\times}5$    \\%
IV & Waited              & \texttt{waited}                           & bin.          & $\beta_k$             & $\beta_k$             & --                    & $\beta_k$             \\%
IV & Reserved            & \texttt{reserved}                         & bin.          & $\beta_k$             & $\beta_k$             & $\beta_k$             & $\beta_k$             \\%
IV & Sentiment           & \texttt{sentiment}                        & cont.         & $\beta_k$             & --                    & --$^b$                & $\beta_k$             \\%
IV & Tag themes          & \texttt{tag\_*} ($\times$10)              & bin.          & --                    & --                    & --                    & --$^c$                \\%
IV & Has text            & \texttt{has\_text}                        & bin.          & --                    & $\beta_k$             & $\beta_k$             & --                    \\%
IV & Visit info          & \texttt{has\_visit\_info}                 & bin.          & $\beta_k$             & $\beta_k$             & --                    & $\beta_k$             \\\hline%
\multicolumn{8}{l}{\quad\textit{Control variables}}\\\hline%
Ctrl & Text Length       & \texttt{ln\_length}                       & cont.         & $\beta_k$             & --                    & $\beta_k$             & $\beta_k$             \\%
Ctrl & User Reviews      & \texttt{ln\_user\_reviews}                & cont.         & $\beta_k$             & $\beta_k$             & $\beta_k$             & $\beta_k$             \\%
Ctrl & User Followers    & \texttt{ln\_user\_followers}              & cont.         & --                    & --                    & --                    & $\beta_k$             \\%
Ctrl & Review Age        & \texttt{ln\_review\_age}                  & cont.         & --                    & --                    & --                    & $\gamma$              \\%
Ctrl & Season            & \texttt{season\_*} ($\times$3)            & bin.          & --                    & --                    & --                    & $\beta_k{\times}3$    \\%
Ctrl & Weekend           & \texttt{visit\_weekend}                   & bin.          & --                    & --                    & --                    & $\beta_k$             \\%
Ctrl & Verification      & \texttt{verif\_receipt}                   & bin.          & --                    & --                    & --                    & $\beta_k$             \\%
Ctrl & Place FE          & \textit{(absorbed)}                       & categ.        & $\delta_j$            & $\delta_j$            & $\delta_j$            & $\delta_j$            \\\hline%
\end{tabular}%
\par\vspace{3pt}\footnotesize%
\textit{Scale:} bin.\ = binary (0/1); ord.\ = ordinal; cont.\ = continuous (log-transformed); count = non-negative integer.\quad
\textit{Notation link:}
$\beta_k$ = enters $\mathbf{x}_{ij}^\top\boldsymbol\beta$ with coefficient $\beta_k$;\quad
$Y_{ij}$ = outcome (left-hand side; not in $\mathbf{x}_{ij}$);\quad
$\delta_j$ = place fixed effect;\quad
$\gamma$ = \texttt{ln\_review\_age} as exposure offset in M4 (not a $\beta_k$);\quad
$\beta_k{\times}N$ = $N$ dummies each with their own $\beta_k$.\\%
$^a$~Partial dummies in M3: only 32\% of reviews have a rating.~%
$^b$~Excluded from M3: $r{=}0.51$ with rating; VIF risk.~%
$^c$~Tag themes only in extended models M1b, M1x.%
\end{table}%
}

\newcommand{\tblVariableMapKO}{%
\begin{table}[H]\centering%
\caption{변수 매핑표}%
\footnotesize\setlength{\tabcolsep}{3.5pt}%
\begin{tabular}{lp{2.7cm}p{2.8cm}lp{1.4cm}p{1.4cm}p{1.4cm}p{1.4cm}}\hline%
유형 & 변수 이름 (표~2) & 데이터셋 열 & 척도 & M1 & M2 & M3 & M4 \\\hline%
\multicolumn{8}{l}{\quad\textit{종속변수}}\\\hline%
종속 & 재방문              & \texttt{revisit}                          & 이진          & $Y_{ij}$              & $\beta_k$             & $\beta_k$             & $\beta_k$             \\%
종속 & 키워드 수           & \texttt{n\_tags}                          & 순서 0--5     & $\beta_k$             & $Y_{ij}$              & $\beta_k$             & $\beta_k$             \\%
종속 & 별점               & \texttt{rating}                           & 순서 1--5     & --                    & --                    & $Y_{ij}$              & --                    \\%
종속 & 반응 수            & \texttt{reactions}                        & 계수          & --                    & --                    & --                    & $Y_{ij}$              \\\hline%
\multicolumn{8}{l}{\quad\textit{독립변수}}\\\hline%
독립 & 동반자 유형         & \texttt{comp\_*} ($\times$6)              & 이진          & $\beta_k{\times}6$    & $\beta_k{\times}6$    & $\beta_k{\times}4^a$  & $\beta_k{\times}6$    \\%
독립 & 방문 목적           & \texttt{purp\_*} ($\times$5)              & 이진          & $\beta_k{\times}5$    & $\beta_k{\times}5$    & $\beta_k{\times}2^a$  & $\beta_k{\times}5$    \\%
독립 & 대기 경험           & \texttt{waited}                           & 이진          & $\beta_k$             & $\beta_k$             & --                    & $\beta_k$             \\%
독립 & 사전 예약           & \texttt{reserved}                         & 이진          & $\beta_k$             & $\beta_k$             & $\beta_k$             & $\beta_k$             \\%
독립 & 감성 점수           & \texttt{sentiment}                        & 연속          & $\beta_k$             & --                    & --$^b$                & $\beta_k$             \\%
독립 & 키워드 주제         & \texttt{tag\_*} ($\times$10)              & 이진          & --                    & --                    & --                    & --$^c$                \\%
독립 & 텍스트 입력 여부    & \texttt{has\_text}                        & 이진          & --                    & $\beta_k$             & $\beta_k$             & --                    \\%
독립 & 방문 정보 제공      & \texttt{has\_visit\_info}                 & 이진          & $\beta_k$             & $\beta_k$             & --                    & $\beta_k$             \\\hline%
\multicolumn{8}{l}{\quad\textit{통제변수}}\\\hline%
통제 & 텍스트 길이         & \texttt{ln\_length}                       & 연속          & $\beta_k$             & --                    & $\beta_k$             & $\beta_k$             \\%
통제 & 작성자 리뷰 수      & \texttt{ln\_user\_reviews}                & 연속          & $\beta_k$             & $\beta_k$             & $\beta_k$             & $\beta_k$             \\%
통제 & 팔로워 수           & \texttt{ln\_user\_followers}              & 연속          & --                    & --                    & --                    & $\beta_k$             \\%
통제 & 리뷰 경과일         & \texttt{ln\_review\_age}                  & 연속          & --                    & --                    & --                    & $\gamma$              \\%
통제 & 방문 계절           & \texttt{season\_*} ($\times$3)            & 이진          & --                    & --                    & --                    & $\beta_k{\times}3$    \\%
통제 & 주말 방문           & \texttt{visit\_weekend}                   & 이진          & --                    & --                    & --                    & $\beta_k$             \\%
통제 & 인증 (영수증)       & \texttt{verif\_receipt}                   & 이진          & --                    & --                    & --                    & $\beta_k$             \\%
통제 & 관광지 고정효과     & \textit{(흡수됨)}                         & 범주          & $\delta_j$            & $\delta_j$            & $\delta_j$            & $\delta_j$            \\\hline%
\end{tabular}%
\par\vspace{3pt}\footnotesize%
\textit{척도:} 이진 = 0/1 더미; 순서 = 순서형 정수; 연속 = 로그 변환 연속형; 계수 = 비음 정수.\quad
\textit{기호 연결:}
$\beta_k$ = $\mathbf{x}_{ij}^\top\boldsymbol\beta$ 벡터에 계수 $\beta_k$로 포함;\quad
$Y_{ij}$ = 종속변수(좌변; $\mathbf{x}_{ij}$에 미포함);\quad
$\delta_j$ = 관광지 고정효과;\quad
$\gamma$ = M4에서 노출 오프셋($\beta_k$ 아님);\quad
$\beta_k{\times}N$ = 더미 $N$개가 각각 고유한 계수로 포함.\\%
$^a$~M3 일부 더미: 별점 보유 리뷰 32\%에 불과.~%
$^b$~M3 제외: 별점과 $r{=}0.51$, 다중공선성 우려.~%
$^c$~키워드 주제는 확장 모형 M1b, M1x에서만.%
\end{table}%
}


\title{}
\date{}

\begin{document}
\pagenumbering{arabic}

% ═══════════════════════════════════════════════════════════════════════
% ENGLISH
% ═══════════════════════════════════════════════════════════════════════

\begin{center}
{\LARGE\bfseries What Makes an Attraction Review Informative}\\[4pt]
{\Large Visit Context, Information Diversity and Reader Response on Naver Place}\\[10pt]
{\normalsize Jorge Adan Soriano Hernandez}\\
{\small Department of Data Science, Hoseo University}\\
{\small Research report, October 2026}
\end{center}

\vspace{6pt}
\subsect{Abstract}
Online reviews are the information travellers consult before deciding whether
to visit or revisit a place; this makes the \emph{diversity} of what a review
says, not just its volume, a determinant of a destination's performance. Lee,
Jeong and Kim (2026) establish this at the business--week level on Naver
Smart Place data. This report takes the same construct to the level of the
individual review: using 14,401 Naver Place reviews of 29 Korean attractions,
it asks which visit contexts predict a reviewer returning, which contexts and
reviewers produce more diverse reviews, and which reviews other users
actually respond to. Four review-level models with place fixed effects
(revisit, keyword diversity, star rating, reactions) recover a consistent
pattern: information diversity tracks a reviewer's engagement with the
platform rather than the nature of the visit, and more positive reviews
receive fewer reactions, replicating at the review level the negative
sentiment effect Lee et al.\ (2026) report on clicks.

\textbf{Keywords:} information diversity, online reviews, review helpfulness,
revisit intention, Naver Place

\sect{I. Introduction}

\subsect{1.1 Research Background}
When deciding whether to visit or revisit a place, one of the strongest
considerations travellers weigh is checking its reviews, comments,
reactions and photos, a direct expression of social media's role in
decision-making. Online reviews are an accessible information source,
unconstrained by time or space, that plays an important role in product and
destination selection (Dellarocas, 2003, as cited in Lee et al., 2026).

South Korea's inbound tourism has grown substantially in recent years, a
trend research attributes partly to the Korean Wave (Bae, Chang, Park \&
Kim, 2017) and to specific experiential segments such as food tourism
(Chang, 2021; Ko, Kang, Kang \& Lee, 2018). Within this context, information
diversity, that is, how distinct the words across a set of reviews actually
are, is understood to drive a place's performance independently of review volume
(Lee, Jeong \& Kim, 2026). This report extends that question to multiple
points of interest across Korea, analysing individual Naver Place reviews
and the diversity of their content alongside the visit context in which they
were written.

Throughout this paper, we operationalize \textbf{visit context} as four
dimensions Naver asks reviewers to report when posting a review: (i)
\textit{companion type} (six categories: alone, partner/spouse, friends,
children, family, colleagues), (ii) \textit{visit purpose} (five categories:
travel, outing, daily life, date, socializing), (iii) \textit{wait experience}
(binary: did the reviewer queue before entry), and (iv) \textit{advance
reservation} (binary). These four dimensions are platform-structured
metadata that few other review ecosystems record, and together they
constitute the central set of predictors examined throughout this study.

\todo{A general tourism/consumer-behaviour citation for "online reviews as a
form of social proof in travel decisions" beyond Dellarocas (2003) would
strengthen this paragraph. Search: "electronic word of mouth travel decision
making" or "social proof online reviews tourism".}

\subsect{1.2 Research Questions and Purpose}
A place's performance can be measured in several ways, yet reviews
themselves rarely separate \emph{how much} was written from \emph{how
informative} it was. Lee et al.\ (2026) show that information diversity
predicts business performance independently of review quantity, but measure
performance through platform-side traffic (page views, clicks) aggregated by
business and week, leaving the review itself, and the context in which it
was written, unobserved. A review written after queueing, with a partner,
on a return visit is treated identically to one written alone, on a first
visit, with no wait, although Naver records all of this. This report asks:

\begin{enumerate}[leftmargin=16pt, itemsep=2pt]
  \item Which visit contexts are associated with returning to the same
        attraction?
  \item Which reviewers and visit contexts produce more diverse information,
        i.e.\ more keyword tags?
  \item Which reviews do other users respond to, and does sentiment help or
        hurt?
  \item Which aspects of the visit (historical, accessibility, crowding,
        food, visual) do different reviewer groups actually talk about, and
        does the content of a review (what is mentioned) predict reader
        engagement beyond its overall sentiment?
\end{enumerate}

\subsect{1.3 Contribution}
This report positions itself with respect to two reference studies and
one methodological absence in the review-analytics literature. Lee,
Jeong and Kim (2026) measure information diversity by aggregating
keyword variety and sentiment to the business-week level, and link it to
platform-side traffic (Page Views, Total Clicks); they do not observe
the review itself or its author. Wang, Karmakar, Wang and Pei (2025)
measure review \emph{helpfulness} as the content dissimilarity of a
review to the rest of the corpus, treating a review as a single bag of
words. Neither study records the visit context in which a review was
written, nor the aspects it covers.

This report makes three contributions on top of that base.
\emph{(i) Review-level replication.} We re-estimate the negative
sentiment-to-engagement effect Lee et al.\ (2026) report at the
business-week level directly at the review level (Model~4,
\S4.4), providing the first review-level corroboration of their finding
in a comparable platform.
\emph{(ii) Platform-structured visit context as a predictor set.} The
four Naver-reported dimensions defined in \S1.1 -- companion type,
visit purpose, wait experience, and reservation -- are platform-captured
metadata that few other review ecosystems record; this is the first
study to use them as a joint predictor set across revisit intention,
information diversity, and reader engagement (Models~1, 2, 4).
\emph{(iii) Aspect-level analysis cross-cut by visit context.} In
\S4.6 we combine LLM-extracted aspect mentions with the visit-context
dimensions and ask who mentions what, and whether the content a review
discusses -- not just how positive it is -- predicts reader engagement.
This cross-cut is novel because it requires both aspect labels and
structured visit metadata on the same reviews, a combination the
reference studies cannot build.

\sect{II. Theoretical Background}

\subsect{2.1 Information Diversity and Business Performance}
Lee, Jeong and Kim (2026) distinguish information \textbf{diversity} --
the number of distinct, non-repeated informational items a review conveys --
from information \textbf{quantity}, the raw amount of content. Working with
Naver Smart Place data aggregated by business and week ($N = 2{,}662{,}298$
business-week observations), they model two outcomes, Page View and Total
Clicks, as functions of Unique Keyword (diversity), Keyword Count
(quantity), review Volume, Text Length and Sentiment, interacted with a
dummy for Naver's March 2023 keyword-expansion update (Expand), under a
two-way (business and week) fixed-effects panel regression. Unique Keyword
has no effect on Page View but a positive effect on Total Clicks; Volume and
Text Length raise both outcomes, with Text Length's effect reversing sign
after the keyword expansion; Sentiment is negative on both outcomes before
the expansion and positive after. The present report keeps their own 0--5
Unique Keyword scale, since Naver caps keyword selection at five, and asks
what produces it at the level of the individual review, rather than what it
produces at the level of the week.

\subsect{2.2 Review Content and Helpfulness}
A directly related strand of research asks what makes a review itself
\emph{useful} to other readers, rather than what makes an aggregate of
reviews move platform traffic. Wang, Karmakar, Wang and Pei (2025) study how
the dissimilarity of a review's content from its surrounding context, namely
product descriptions and preceding reviews, shapes its perceived
helpfulness, using roughly 1.6 million Amazon reviews. They report that
findings in this area have been mixed: Fresneda and Gefen (2019, as cited in
Wang et al., 2025) find a \emph{positive} relationship between the count of
unique words in a review and its helpfulness, while Namvar and Chua (2023,
as cited in Wang et al., 2025) find a \emph{negative} effect of a review's
lexical dissimilarity from preceding reviews on helpfulness. This divergence
is part of the motivation for treating the reactions a review receives
(Model 4, below) as its own outcome rather than assuming diversity and
helpfulness move together.

\subsect{2.3 Review Sentiment and Reader Response}
Sentiment is scored in this report with a BERT-family Korean sentiment
classifier (hun3359/klue-bert-base-sentiment); BERT-based sentiment models
have been shown to substantially outperform earlier approaches in
consumer-review settings, reaching 91\% accuracy in product-recommendation
contexts (Karabila, Darraz, EL-Ansari, Alami \& EL Mallahi, 2024). Beyond the
classifier itself, the theoretical question of \emph{why} a more positive
review might attract fewer reader reactions, the pattern both this report
(Section IV) and Lee et al.\ (2026) observe, is not addressed by a source
read for this project.

\todo{Theoretical grounding for why positive sentiment can reduce reader
engagement (e.g.\ negativity bias in word-of-mouth, or diagnosticity of
negative information) would strengthen this subsection. Search: "negativity
bias online reviews helpfulness" or "eWOM valence reader engagement".}

\subsect{2.4 Research Gap}
Taken together, the literature establishes that information diversity
matters for performance (Lee et al., 2026) and that review content shapes
perceived helpfulness (Wang et al., 2025), but neither source observes the
\emph{antecedents} of diversity at the review level: which visitors, under
which circumstances, actually produce an informative review. Naver records
the visit's companion, purpose and waiting time for every review; this
report is, to the author's knowledge, the first to connect that visit
context to information diversity, revisit behaviour and reader response in
a single review-level framework.

\sect{III. Methodology}

\subsect{3.1 Data}
Naver Place reviews for Korean attractions, collected by scraping each place
individually. The dataset holds 14,401 reviews from 8,276 reviewers across
29 attractions; 11,165 carry review text and 4,651 carry a star rating. The
unit of analysis is the review $i$ at place $j$, not the business-week.
Places were selected by crossing reputational prominence (Yanolja Research
"Korea's Top 500") with visitor volume by origin (Korea Tourism
Organization), inside typological strata (palaces, traditional quarters,
markets, entertainment, museums, modern landmarks, nature).

\subsect{3.2 Variables}
Four dependent variables, the visit-context regressors Naver records, and
the controls that separate a behavioural effect from an artefact of who
wrote the review and when.

\tblVariables

The choice of independent variables rests on three strands of prior work.
\textbf{Companion type and visit purpose.} Tourism segmentation research
treats \emph{who you travel with} and \emph{why you came} as the two
dimensions that most strongly shape what a visitor experiences and how
they evaluate it afterwards. Ko, Kang, Kang and Lee (2018) use a
factor-cluster segmentation on visit purpose to classify foreign tourists'
perceptions of Korean food tours into interpretable groups; Chang (2021)
treats food tourism itself as a distinct experiential segment with its own
antecedents and outcomes; Bae, Chang, Park and Kim (2017) tie segment
growth to the Korean Wave. If the kind of visit shapes what a visitor
brings home to say about it, these two Naver-reported fields should enter
review-level models of what is written and how it is received, so we use
the full six-category companion dummies and five-category purpose dummies.
\textbf{Wait experience and reservation.} These are operational signals
of the service encounter, not sentiment: queueing is the most salient
friction in a visitor's assessment, and prior reservation changes both
baseline expectation and the perceived stakes of the visit. We treat both
as binary descriptors of the encounter and let the regressions discover
whether they map onto review behaviour; \S4.6 confirms that they do.
\textbf{Sentiment, keyword themes and reader response.} Sentiment is scored
with the BERT-based Korean classifier described in \S2.3 (Karabila et al.,
2024); Keyword Count operationalizes information diversity in the exact
0--5 band Lee et al.\ (2026) use; and the reactions a Naver review
accumulates serve as this report's behavioural proxy for review helpfulness
(Wang et al., 2025), the dimension \S2.2 motivates.

\subsect{3.3 Descriptive Statistics}
\tblDescriptives

Keyword Count sits on the same 0--5 scale Lee et al.\ (2026) report, because
Naver caps keyword selection at five. The star rating moves little (median
5), which is why Model 3 is reported as exploratory rather than a primary
result.

\subsect{3.4 Correlation Analysis}
\tblCorrelations

No pair threatens the regressions below. The largest, Rating with Sentiment
(0.514), is the evidence that the sentiment classifier tracks the stars
reviewers actually gave, not a collinearity concern. Keyword Count correlates
$-0.395$ with Review Age: older reviews carry fewer keywords, so keyword use
has grown over the observation window, the same keyword-expansion dynamic
Lee et al.\ (2026) exploit as an event in their own design.

\subsect{3.5 Models}
Review-level regressions with place fixed effects $\delta_j$ and standard
errors clustered by reviewer, since the same person often writes several
reviews. The link function follows the measurement scale of each outcome:
logistic for binary Revisit (M1), ordered logistic for Keyword Count (M2)
and Rating (M3), and negative binomial for the overdispersed Reactions
count (M4). Keyword Count is modeled as ordinal rather than as a count
because its range is capped by Naver's platform design at five tags, not
generated by a Poisson-distributed data process; the same reasoning
applies to the 1--5 star Rating.

\vspace{-4pt}
\begin{align*}
  \text{M1 (Revisit):}\quad &
    \operatorname{logit} P(\textit{Revisit}_{ij}) = \beta_0 + \mathbf{x}_{ij}^\top\boldsymbol\beta + \delta_j \\[2pt]
  \text{M2 (Keyword Count):}\quad &
    P(\textit{KeywordCount}_{ij} \le k) = \big[1 + e^{-(\tau_k - \mathbf{x}_{ij}^\top\boldsymbol\beta - \delta_j)}\big]^{-1},\ k = 0,\dots,4 \\[2pt]
  \text{M3 (Rating):}\quad &
    P(\textit{Rating}_{ij} \le k) = \big[1 + e^{-(\tau_k - \mathbf{x}_{ij}^\top\boldsymbol\beta - \delta_j)}\big]^{-1},\ k = 1,\dots,4 \\[2pt]
  \text{M4 (Reactions):}\quad &
    \textit{Reactions}_{ij} \sim \text{NegBin}(\mu_{ij}),\quad
    \ln \mu_{ij} = \beta_0 + \mathbf{x}_{ij}^\top\boldsymbol\beta + \gamma\ln(\textit{Age}_{ij}) + \delta_j
\end{align*}
\vspace{-6pt}

M4 uses a negative binomial rather than Poisson link because the variance of
Reactions is roughly 37 times its mean in this sample; a Poisson model was
estimated first and rejected on that basis, and the negative binomial's
dispersion parameter is estimated at $\alpha = 2.44$. Each primary model is
re-estimated on one random review per reviewer and per place (the "one per
reviewer$\times$place" sensitivity check); only effects that survive this
check in both the full sample and the sensitivity sample are reported as
robust in Section IV. Out-of-sample metrics (AUC, RMSE) come from 5-fold
cross-validation grouped by reviewer, so no reviewer appears in both the
training and the test fold.

\tblVIF

\sect{IV. Results}

\subsect{4.1 Revisit (Model 1)}
Model 1 asks which visit contexts distinguish a review written on a return
visit from one written on a first visit. Out-of-sample AUC is 0.773 (0.785
for the extended specification, Model 1x, which adds season, weekday,
verification method and reviewer-follower controls; AIC improves from 9{,}596
to 9{,}372).

\begin{table}[H]\centering
\caption{Model 1: P(Revisit), logistic regression, place fixed effects\\
\footnotesize $N = 11{,}165$ reviews with text (6{,}070 reviewers), revisit rate $= 22.6\%$, pseudo-$R^2 = 0.203$}
\footnotesize
\begin{tabular}{@{}lrrr@{}}
\toprule
Variable & OR & 95\% CI & $p$ \\
\midrule
Waited to enter          & 0.51 & [0.35, 0.72] & 0.0002 *** \\
With: alone               & 1.68 & [1.18, 2.39] & 0.0037 ** \\
With: partner/spouse      & 0.74 & [0.53, 1.04] & 0.0807 . \\
With: friends              & 0.58 & [0.42, 0.80] & 0.0011 ** \\
Purpose: travel            & 0.56 & [0.42, 0.75] & 0.0001 *** \\
Purpose: date               & 0.72 & [0.53, 0.97] & 0.0284 * \\
Sentiment ($-1$ to $1$)    & 0.70 & [0.60, 0.81] & 0.0000 *** \\
Keyword tags (0--5)        & 1.14 & [1.07, 1.23] & 0.0002 *** \\
ln(text length)             & 1.29 & [1.14, 1.47] & 0.0001 *** \\
ln(1 + reviewer's reviews) & 1.52 & [1.40, 1.65] & 0.0000 *** \\
\bottomrule
\end{tabular}
\\[3pt]
{\scriptsize Reservation, visit-info-provided and the remaining companion/purpose dummies were not significant at $p<.05$. Full output: \texttt{data/analysis/naver/model1\_revisit.csv}.}
\end{table}

\textbf{Robust across the sensitivity check} (same conclusion at $p<.05$ on
one review per reviewer$\times$place): waited to enter, with: alone, with:
friends, sentiment, keyword tags, ln(1 + reviewer's reviews). \textbf{Not
robust}: visit info provided, with: partner/spouse, purpose: date, ln(text
length) -- these lose significance once repeated reviewers are downweighted
to one review each, and should be read with caution.

In the extended specification (Model 1x), the single largest effect in the
entire table is \textbf{verification by receipt} (OR $= 0.35$, $p<.001$): a
receipt-verified review is far less likely to come from a returning visitor.
ln(1 + reviewer's followers) is also positive (OR $=1.18$, $p<.001$), and
tags mentioning water leisure are associated with lower revisit odds (OR
$=0.62$, $p=.009$).

\subsect{4.2 Keyword Diversity (Model 2)}
Model 2 asks which reviewers and visit contexts produce more diverse
reviews. Out-of-sample RMSE is 1.366 against a mean-prediction baseline of
1.883 (1.337 for the extended specification, Model 2x).

\begin{table}[H]\centering
\caption{Model 2: Keyword Count (0--5), ordered logistic regression, place fixed effects\\
\footnotesize $N = 14{,}401$, mean $= 2.97$ tags}
\footnotesize
\begin{tabular}{@{}lrrr@{}}
\toprule
Variable & OR & 95\% CI & $p$ \\
\midrule
Wrote text (15+ chars)     & 4.02 & [3.51, 4.60] & 0.0000 *** \\
Visit info provided         & 2.53 & [2.01, 3.18] & 0.0000 *** \\
Revisitor (2+ visits)        & 1.60 & [1.32, 1.93] & 0.0000 *** \\
With: alone                 & 1.67 & [1.27, 2.20] & 0.0003 *** \\
Purpose: outing              & 1.85 & [1.60, 2.13] & 0.0000 *** \\
Purpose: travel               & 1.73 & [1.49, 2.02] & 0.0000 *** \\
Purpose: socializing          & 1.59 & [1.31, 1.93] & 0.0000 *** \\
Waited to enter               & 0.75 & [0.64, 0.88] & 0.0004 *** \\
ln(1 + reviewer's reviews)  & 1.17 & [1.13, 1.22] & 0.0000 *** \\
\bottomrule
\end{tabular}
\\[3pt]
{\scriptsize All companion and purpose dummies are directionally positive; only those above $p<.001$ are shown for brevity. Full output: \texttt{data/analysis/naver/model2\_keywords.csv}.}
\end{table}

\textbf{Robust across specifications} (significant in both the ordered-logit
and the binomial robustness check, Model 2b): revisitor, wrote text, visit
info provided, waited to enter, with: alone, with: partner/spouse, with:
parents/relatives, purpose: travel, outing, daily life, date, socializing,
ln(1 + reviewer's reviews). Among reviews with text (Model 2t), sentiment
enters \emph{positively} (OR $=2.26$, $p<.001$) and ln(text length) strongly
so (OR $=1.86$, $p<.001$) -- unsurprising, since longer, more positive text
naturally accompanies more keyword tags, but it is why sentiment and length
are excluded from the main specification (Model 2): they are products of the
same act of writing as the tags themselves, not antecedents of it.

\subsect{4.3 Star Rating (Model 3): Exploratory}
Star ratings move little in this corpus (median 5, 73.9\% five-star), so
Model 3 is reported as exploratory. McFadden $R^2 = 0.095$ ($N = 4{,}651$).
The one effect that is both large and consistent across the base, binomial
(Model 3b) and text-only (Model 3t) specifications is \textbf{Keyword Count}
(OR $= 1.56$, $p<.001$ in the base model): reviewers who tag more keywords
also rate higher. In the text-only specification, Sentiment dominates
everything else (OR $= 4.90$, $p<.001$), as expected given it is scored from
the same text the rating accompanies -- it is a validation of the sentiment
classifier, not a behavioural finding. Wrote text is negatively associated
with the star rating (OR $=0.64$, $p<.001$) -- reviewers who write text rate
slightly lower than those who leave only keyword tags.

\subsect{4.4 Reactions (Model 4)}
Model 4 asks which reviews other users respond to. Following the framing in
\S2.2, we treat the reactions a review accumulates as a behavioural proxy
for review helpfulness (Wang et al., 2025): a review that enough other
viewers found worth clicking on. Out-of-sample AUC for predicting whether
a review receives any reaction is 0.815 from the count model's predicted
mean (0.827 for the dedicated binary specification, Model 4b);
pseudo-$R^2 = 0.212$.

\begin{table}[H]\centering
\caption{Model 4: Reactions per review, negative binomial regression, place fixed effects\\
\footnotesize $N = 11{,}165$ reviews with text, reviews with $\geq 1$ reaction $= 23.5\%$, mean $= 1.82$, $\alpha = 2.44$, pseudo-$R^2 = 0.212$}
\footnotesize
\begin{tabular}{@{}lrrr@{}}
\toprule
Variable & IRR & 95\% CI & $p$ \\
\midrule
ln(1 + reviewer's followers) & 1.97 & [1.79, 2.18] & 0.0000 *** \\
ln(text length)               & 1.64 & [1.44, 1.87] & 0.0000 *** \\
Keyword tags (0--5)            & 1.12 & [1.05, 1.20] & 0.0011 ** \\
\textbf{Sentiment ($-1$ to $1$)} & \textbf{0.77} & [0.64, 0.91] & 0.0032 ** \\
Visit info provided             & 1.70 & [1.08, 2.68] & 0.0223 * \\
Purpose: outing                  & 0.63 & [0.45, 0.88] & 0.0070 ** \\
Purpose: travel                    & 0.62 & [0.46, 0.83] & 0.0016 ** \\
Visit: autumn                        & 0.71 & [0.61, 0.83] & 0.0000 *** \\
ln(1 + reviewer's reviews)         & 0.90 & [0.81, 0.99] & 0.0378 * \\
\bottomrule
\end{tabular}
\\[3pt]
{\scriptsize Companion dummies are not significant at $p<.05$. Full output: \texttt{data/analysis/naver/model4\_reactions.csv}.}
\end{table}

\textbf{This replicates Lee et al.\ (2026)}, who also find sentiment entering
engagement negatively before Naver's keyword expansion ($\alpha_{\text{Sentiment}}
= -0.011$ to $-0.018$ on Page View/Total Clicks). At the review level, more
positive reviews receive measurably fewer reactions -- effusive praise is not
what readers respond to. \textbf{Robust across specifications} (significant
in both the negative binomial and the binary any-reaction model, Model 4b):
keyword tags, ln(text length), visit: autumn, ln(1 + reviewer's followers).
Sentiment and ln(1 + reviewer's reviews) are significant in Model 4 but not
in Model 4b, and should be read as suggestive rather than fully robust.

\subsect{4.5 Model Performance Summary}
\begin{table}[H]\centering
\caption{Out-of-sample performance, 5-fold cross-validation grouped by reviewer}
\footnotesize
\begin{tabular}{@{}llll@{}}
\toprule
Model & Metric & Value & Baseline \\
\midrule
M1 Revisit              & AUC          & 0.773 (ext.\ 0.785) & -- \\
M2 Keyword Count         & RMSE          & 1.366 (ext.\ 1.337)  & 1.883 (predict mean) \\
M3 Rating                & --            & exploratory          & -- \\
M4 Reactions              & AUC (any)      & 0.815                & -- \\
M4b Any reaction           & AUC            & 0.827                & -- \\
\bottomrule
\end{tabular}
\end{table}

\subsect{4.6 Aspect-Level Analysis}
The sentiment score summarizes a review on a single axis, but a review's
content lives on several axes at once. Using an LLM-based classifier
(\texttt{06\_llm\_analysis.py}), we labelled each of the 11{,}164
text-bearing Korean reviews with aspect-level polarity ($-1$/$0$/$+1$)
on six dimensions: \emph{historical/cultural} (50.4\% mentioned),
\emph{visual/photo} (63.8\%), \emph{service/experience} (27.4\%),
\emph{accessibility} (19.8\%), \emph{crowding} (16.3\%) and
\emph{food} (6.9\%).

The six aspects were derived from tourism-literature grounds (SERVQUAL's
tangibles and responsiveness, Parasuraman, Zeithaml \& Berry, 1988; food
tourism as a distinct experiential segment, Chang, 2021; Ko et al., 2018)
and corroborated by unsupervised topic extraction on the same corpus
\emph{prior to} LLM classification: Latent Dirichlet Allocation at
optimal $K=6$ (coherence $c_v = 0.520$) and BERTopic with multilingual
embeddings ($K=15$). All six aspects except crowding emerge as distinct
topics in both unsupervised methods; crowding is retained because it is
a polarity-dominated dimension (18.1\% of its mentions are negative, by
far the highest among aspects) rather than a word-frequency topic. Full
validation details: \texttt{data/analysis/naver/topic\_validation/}.

This subsection asks two questions the main specifications cannot:
\emph{(a)} who mentions which aspect, and \emph{(b)} does what people
mention predict engagement beyond overall sentiment?

\textbf{4.6.1 Who mentions what?} Six parallel logistic regressions
predict whether each aspect is mentioned, using the four visit-context
dimensions, companion/purpose dummies, reviewer controls, and place
fixed effects (standard errors clustered by reviewer). Table~\ref{tab:absaA}
reports the significant visit-context effects (full output:
\texttt{data/analysis/naver/absa\_mention\_*.csv}).

\begin{table}[H]\centering
\caption{Model ABSA-A: Aspect mention as a function of visit context (odds ratios; $N=11{,}164$). Blank cells = $p>.05$.}
\label{tab:absaA}
\footnotesize
\setlength{\tabcolsep}{3pt}
\begin{tabular}{@{}lcccccc@{}}
\toprule
Visit context                  & Hist.\     & Access.\   & Crowd.\    & Food       & Visual     & Service    \\
\midrule
With: alone                     & --         & --         & --         & 0.64 *      & --         & 0.52 **     \\
With: partner/spouse            & --         & --         & 1.25 *      & --         & 1.21 *      & --         \\
With: friends                   & --         & --         & --         & --         & 1.28 *      & --         \\
With: children                  & 1.22 *      & 1.22 *      & --         & --         & --         & 1.55 ***    \\
With: parents/relatives         & --         & --         & --         & --         & 1.28 **     & --         \\
With: colleagues                & --         & --         & 0.72 *      & --         & --         & --         \\
Purpose: outing                 & --         & --         & --         & --         & --         & 0.75 **     \\
Purpose: date                   & --         & --         & 0.76 *      & 1.51 *      & --         & --         \\
Waited to enter                 & --         & --         & 2.13 ***    & --         & --         & --         \\
Reservation                     & --         & --         & 0.69 **     & --         & --         & 2.39 ***    \\
\midrule
ln(text length)                 & 2.18 ***    & 1.67 ***    & 1.85 ***    & 1.40 ***    & 1.67 ***    & 2.30 ***    \\
ln(1 + reviewer's reviews)      & --         & --         & 1.05 *      & --         & 1.13 ***    & 0.79 ***    \\
\midrule
Mention rate                    & 50.4\%      & 19.8\%      & 16.3\%      & 6.9\%       & 63.8\%      & 27.4\%      \\
Pseudo-$R^2$                    & 0.471       & 0.079       & 0.113       & 0.133       & 0.217       & 0.371       \\
\bottomrule
\end{tabular}
\\[3pt]
{\scriptsize Place fixed effects and remaining non-significant companion/purpose dummies not shown. Significance: \textit{***} $p<.001$, \textit{**} $p<.01$, \textit{*} $p<.05$.}
\end{table}

Four findings stand out, all new to this literature.
\textbf{First, children anchor narrative-rich reviews across aspects.}
Visits with children raise the odds of mentioning historical/cultural
content by 22\% ($p=.031$), accessibility by 22\% ($p=.035$), and
service/experience by 55\% ($p<.001$). The pattern is coherent:
caregivers notice and record more of the visit's dimensions than other
visitor types.
\textbf{Second, service friction leaks into content, two ways.}
Queueing at the entrance more than doubles the odds of a crowding
mention (OR $2.13$, $p<.001$), while having a reservation cuts them to
69\% (OR $0.69$, $p=.001$). On the other hand, reservations more than
double the odds of a service/experience mention (OR $2.39$, $p<.001$):
if you booked ahead, you came expecting service and you comment on it.
\textbf{Third, dating anchors food talk.} Date-purpose visits raise the
odds of food mentions by 51\% (OR $1.51$, $p=.028$), while alone visits
lower them (OR $0.64$, $p=.026$) -- consistent with food's social role
in Korean dating culture (Chang, 2021).
\textbf{Fourth, service/experience is a distinct dimension in the
Korean platform.} The service aspect has the second-highest pseudo-$R^2$
(0.371, after historical/cultural) and its predictors are structurally
different from any other aspect: strongly positive for reservations and
children, strongly negative for alone visits and higher reviewer
experience. This is the empirical justification for adding it as the
sixth aspect alongside Lee et al.'s (2026) five tourism-standard
dimensions.

\textbf{4.6.2 Do aspect mentions predict engagement?} We add the six
aspect-mention dummies to a reduced Model~4 specification on the ABSA
sample ($N=11{,}164$; a reduced spec drops sparse controls so the
negative binomial converges). Table~\ref{tab:absaB} reports the six
dummies' effects.

\begin{table}[H]\centering
\caption{Model ABSA-B: Reactions as a function of aspect mentions (IRRs; $N=11{,}164$, place FE, log-age offset, SE clustered by reviewer).}
\label{tab:absaB}
\footnotesize
\begin{tabular}{@{}lrrr@{}}
\toprule
Aspect mentioned             & IRR   & 95\% CI               & $p$      \\
\midrule
Historical/cultural          & 0.73  & [0.52,\, 1.00]        & 0.053 .   \\
Accessibility                & 0.89  & [0.72,\, 1.09]        & 0.257    \\
Crowding                     & 1.14  & [0.94,\, 1.39]        & 0.185    \\
Food                         & 0.96  & [0.71,\, 1.31]        & 0.806    \\
Visual/photo                 & 0.84  & [0.63,\, 1.13]        & 0.254    \\
Service/experience           & 1.07  & [0.82,\, 1.40]        & 0.624    \\
\bottomrule
\end{tabular}
\\[3pt]
{\scriptsize Significance: $.\ p<.10$, no coefficient reaches $p<.05$. Full output including non-aspect covariates: \texttt{data/analysis/naver/absa\_reactions\_extended.csv}.}
\end{table}

The pattern is striking: \textbf{no aspect-mention dummy reaches
statistical significance at $p<.05$}, with historical/cultural
approaching marginal significance ($\text{IRR}=0.73$, $p=.053$). The
reactions model's explanatory power stays concentrated in the
reviewer-identity variables already documented in \S4.4: follower count
(IRR $2.29$, $p<.001$) and revisitor status (IRR $2.25$, $p<.001$)
remain the dominant predictors, with keyword tags marginal
($\text{IRR}=1.15$, $p=.051$).

The interpretation is a substantive finding, not an estimation failure:
in Korean tourism reviews, reader engagement is driven by \emph{who
wrote the review} and the review's quality signals (length, keyword use),
not by \emph{what specific aspect} the review happens to discuss. Once
reviewer identity is accounted for, the content-aspect level adds no
measurable explanatory power. This refines the \S4.4 sentiment finding:
sentiment polarity (which varies across reviews) and content aspect
(which varies in what is discussed) do not act on reactions in the same
way -- sentiment polarity still associates with reactions through M4's
compound sentiment variable, but what the review \emph{talks about} does
not. For platform design, this points toward reviewer-identity signals
(follower counts, reviewer history) as the actionable levers for
engagement, rather than nudging reviewers toward particular aspects.

\sect{V. Conclusion and Implications}

\subsect{5.1 Summary of Findings}
\begin{itemize}[leftmargin=14pt, itemsep=2pt]
  \item Information diversity (Keyword Count) on Naver Place tracks how
        engaged a reviewer is with the platform -- whether they wrote text,
        provided visit details, have reviewed often before -- more than it
        tracks the kind of visit they made.
  \item Sentiment enters reader engagement (reactions) negatively,
        replicating Lee et al.\ (2026) at the individual-review level rather
        than the business-week level.
  \item Queueing (waited to enter) is the single strongest negative signal
        for returning to an attraction, and the only one of these findings
        an attraction manager can directly act on.
  \item Star ratings carry comparatively little information in this corpus;
        keyword tags are the more informative dependent variable.
\end{itemize}

\subsect{5.2 Implications}
For platforms, information diversity tracks reviewer engagement rather than
visit type, so prompting for varied keywords should work better aimed at
experienced reviewers than at first-time visitors. For attraction managers,
reducing wait times is the one lever in this report's results that plausibly
affects whether a visitor returns.

\subsect{5.3 Limitations and Next Steps}
\begin{itemize}[leftmargin=14pt, itemsep=2pt]
  \item Associations only. Lee et al.\ (2026) identify effects from Naver's
        keyword-expansion event and lag all regressors to $t-1$; this design
        is cross-sectional and has neither.
  \item Only place fixed effects are included, where Lee et al.\ (2026) use
        place \emph{and} time fixed effects. Month fixed effects, and a
        formal Hausman test to justify fixed over random effects, are the
        next addition.
  \item Model 1 predicts a visit that already happened from a review written
        afterwards, so it profiles returning visitors rather than explaining
        why anyone returns.
  \item A single platform: nothing here separates Naver's reviewer
        population from Korean reviewers in general.
\end{itemize}

\vspace{8pt}
\subsect{References}
\begin{itemize}[leftmargin=14pt, itemsep=3pt]
  \item Bae, E., Chang, M., Park, E., \& Kim, D. (2017). The effect of Hallyu
        on tourism in Korea. \textit{Journal of Open Innovation: Technology,
        Market, and Complexity}, 3, 22.
  \item Chang, S. (2021). Food tourism in Korea. \textit{Journal of Vacation
        Marketing}.
  \item Ko, S., Kang, S., Kang, H., \& Lee, M. J. (2018). An exploration of
        foreign tourists' perceptions of Korean food tour: a factor-cluster
        segmentation approach. \textit{Asia Pacific Journal of Tourism
        Research}, 23(8), 833--846.
  \item Karabila, I., Darraz, N., EL-Ansari, A., Alami, N., \& EL Mallahi, M.
        (2024). BERT-enhanced sentiment analysis for personalized e-commerce
        recommendations. \textit{Multimedia Tools and Applications}, 83,
        56463--56488.
  \item Lee, M., Jeong, D., \& Kim, A. J. (2026). The impact of information
        diversity on online reviews and business performance.
        \textit{Journal of Consumer Studies}, 37(1).
        \url{https://doi.org/10.35736/JCS.37.1.1}
  \item Wang, S., Karmakar, S., Wang, F., \& Pei, Y. (2025). Content
        dissimilarity and online review helpfulness: Contextual insights.
        \textit{Journal of Business Research}, 187, 115068.
\end{itemize}

\vspace{4pt}
{\footnotesize\textit{Fresneda \& Gefen (2019), Namvar \& Chua (2023) and
Dellarocas (2003) are cited above only as reported inside Wang et al.\
(2025) and Lee et al.\ (2026) respectively -- they were not read directly for
this report and should be verified against the original sources before this
report is submitted anywhere. The three \texttt{[SOURCE NEEDED]} markers in
Sections 1.1, 2.3 and 2.4 above name exactly what additional literature
would be needed and are not placeholders for invented citations.}}

\clearpage
\setcounter{table}{0}

% ═══════════════════════════════════════════════════════════════════════
% KOREAN
% ═══════════════════════════════════════════════════════════════════════

\begin{center}
{\LARGE\bfseries 무엇이 관광지 리뷰를 유익하게 만드는가}\\[4pt]
{\Large 네이버 플레이스의 방문 맥락, 정보 다양성, 이용자 반응}\\[10pt]
{\normalsize 소리아노 에르난데스 호르헤 아단}\\
{\small 호서대학교 데이터사이언스학과}\\
{\small 연구 보고서, 2026년 10월}
\end{center}

\vspace{6pt}
\subsect{국문초록}
온라인 리뷰는 여행자가 장소를 방문하거나 재방문할지 결정할 때 참고하는 정보이며,
이로 인해 리뷰의 양뿐 아니라 \textbf{다양성} — 리뷰가 전달하는 내용이 얼마나
다른가 — 이 관광지 성과를 결정하는 요인이 된다. 이민영·정동영·김지영(2026)은
네이버 스마트 플레이스 자료를 업체·주 단위로 집계하여 이를 입증하였다. 본
보고서는 동일한 개념을 개별 리뷰 수준으로 확장하여, 한국 관광지 29곳의 네이버
플레이스 리뷰 14,401건을 활용해 어떤 방문 맥락이 재방문과 연관되는지, 어떤
맥락과 작성자가 더 다양한 리뷰를 생산하는지, 그리고 어떤 리뷰가 다른 이용자의
반응을 얻는지를 묻는다. 관광지 고정효과를 포함한 네 개의 리뷰 단위 모형(재방문,
키워드 다양성, 별점, 반응)은 일관된 패턴을 보여준다. 정보 다양성은 방문 유형보다
작성자의 플랫폼 참여도를 따르며, 더 긍정적인 리뷰일수록 반응은 오히려 적다 —
이는 이민영 외(2026)가 클릭 수에서 보고한 부정적 감성 효과를 개별 리뷰 수준에서
재현한 결과이다.

\textbf{주제어:} 정보 다양성, 온라인 리뷰, 리뷰 유용성, 재방문 의도, 네이버 플레이스

\sect{I. 서론}

\subsect{1.1 연구배경}
장소를 방문하거나 재방문할지 결정할 때, 여행자가 가장 중요하게 고려하는 것 중
하나는 리뷰·댓글·반응·사진을 확인하는 일이며, 이는 소셜 미디어가 의사결정에
미치는 영향을 직접적으로 보여준다. 온라인 리뷰는 시간과 공간의 제약 없이 접근
가능한 정보원으로서, 제품 및 목적지 선택 과정에서 중요한 역할을 한다(Dellarocas,
2003; 이민영 외, 2026에서 재인용).

한국의 인바운드 관광은 최근 크게 성장하였으며, 이는 한류(Bae, Chang, Park \&
Kim, 2017)와 음식관광과 같은 특정 체험 세그먼트(Chang, 2021; Ko, Kang, Kang \&
Lee, 2018)에 부분적으로 기인하는 것으로 연구되었다. 이러한 맥락에서 정보
다양성 — 리뷰에 사용된 단어가 서로 얼마나 다른가 — 은 리뷰 양과 무관하게
장소의 성과를 이끄는 요인으로 이해된다(이민영 외, 2026). 본 보고서는 이 질문을
한국의 여러 관광지로 확장하여, 개별 네이버 플레이스 리뷰와 그 내용의 다양성을
작성 당시의 방문 맥락과 함께 분석한다.

\todo{Dellarocas(2003) 외에 "온라인 리뷰가 여행 의사결정에서 사회적 증거로
작용한다"는 주장을 뒷받침할 일반 관광/소비자행동 문헌이 추가로 필요하다.
검색어: "electronic word of mouth travel decision making" 또는
"social proof online reviews tourism".}

\subsect{1.2 연구 질문 및 목적}
장소의 성과는 여러 방식으로 측정될 수 있지만, 리뷰 자체는 "얼마나 많이
썼는가"와 "얼마나 유익하게 썼는가"를 거의 구분하지 않는다. 이민영 외(2026)는
정보 다양성이 리뷰 양과 무관하게 업체 성과를 예측함을 보였으나, 성과를
업체·주 단위로 집계된 플랫폼 측 트래픽(페이지 뷰, 클릭)으로 측정하여 리뷰
자체와 그 작성 맥락은 관측하지 못한다. 대기 후 연인과 함께 재방문하며 쓴
리뷰와, 혼자 대기 없이 첫 방문에 쓴 리뷰가 동일하게 취급된다 — 네이버가 이
모든 정보를 기록하고 있음에도 말이다. 본 보고서는 다음을 묻는다.

\begin{enumerate}[leftmargin=16pt, itemsep=2pt]
  \item 어떤 방문 맥락이 동일 관광지로의 재방문과 연관되는가?
  \item 어떤 작성자와 방문 맥락이 더 다양한 정보, 즉 더 많은 키워드를
        제공하는가?
  \item 어떤 리뷰가 다른 이용자의 반응을 얻으며, 감성은 이에 도움이 되는가
        아니면 저해하는가?
  \item 어떤 방문 측면을 작성자 그룹별로 언급하며, 리뷰가 다루는
        \emph{내용}이 전반 감성 이상의 이용자 반응을 예측하는가?
\end{enumerate}

\subsect{1.3 연구 기여}
본 연구는 리뷰 분석 문헌의 두 참조 연구와 하나의 방법론적 공백에 대해
자리매김한다. 이민영 외(2026)는 키워드 다양성과 감성을 사업체·주 단위로
집계하여 플랫폼 트래픽(페이지 뷰, 클릭)과 연결하지만, 리뷰 자체나 작성자는
관측하지 않는다. 왕 외(2025)는 리뷰 \emph{유용성}을 콘텐츠 비유사성으로
측정하며, 리뷰를 하나의 단어 집합으로 다룬다. 두 연구 모두 리뷰가 작성된
방문 맥락이나 다루는 측면은 기록하지 않는다.

본 연구는 이 기반 위에 세 가지 기여를 추가한다.
\emph{(i) 리뷰 수준 재현.} 이민영 외(2026)가 사업체·주 단위에서 보고한
부정적 감성-몰입 효과를 리뷰 수준에서 재추정한다(모형 4, \S4.4).
\emph{(ii) 플랫폼 구조화 방문 맥락을 통합 예측 변수 세트로 활용.} \S1.1에서
정의한 네 가지 네이버 수집 차원---동반자 유형, 방문 목적, 대기 경험, 사전
예약---은 다른 리뷰 생태계에서 거의 기록되지 않는 플랫폼 포착 메타데이터이며,
재방문 의도, 정보 다양성, 독자 반응에 걸쳐 통합 예측 변수 세트로 사용한 것은
본 연구가 처음이다(모형 1, 2, 4).
\emph{(iii) 방문 맥락과 교차한 측면 수준 분석.} \S4.6에서 LLM으로 추출한
측면 언급을 방문 맥락 차원과 결합하여 누가 무엇을 언급하는지, 그리고 리뷰
내용이 전반 감성 이상의 반응을 예측하는지 분석한다. 이 교차 분석은 동일
리뷰에 측면 라벨과 구조화된 방문 메타데이터가 모두 있어야 가능하며,
참조 연구들은 구성할 수 없는 조합이다.

\sect{II. 이론적 배경}

\subsect{2.1 정보 다양성과 업체 성과}
이민영·정동영·김지영(2026)은 리뷰가 전달하는 중복되지 않는 정보 항목의 수인
정보 \textbf{다양성}을, 콘텐츠의 총량인 정보 \textbf{양}과 구분한다. 업체·주
단위로 집계된 네이버 스마트 플레이스 자료($N = 2{,}662{,}298$ 업체-주
관측값)를 활용하여, 페이지 뷰와 총 클릭 수를 Unique Keyword(다양성), Keyword
Count(양), 리뷰 수(Volume), 텍스트 길이, 감성 점수의 함수로, 2023년 3월
네이버의 키워드 확장 업데이트 전후 더미(Expand)와의 상호작용을 포함하여,
양방향(업체·주) 고정효과 패널 회귀분석으로 추정하였다. Unique Keyword는
페이지 뷰에는 유의한 영향이 없으나 총 클릭 수에는 긍정적인 영향을 미치며,
Volume과 Text Length는 두 종속변수 모두를 증가시키되 Text Length의 효과는
키워드 확장 이후 방향이 전환된다. Sentiment는 확장 이전에는 두 종속변수
모두에 부정적이나 확장 이후에는 긍정적으로 전환된다. 본 보고서는 이들의 0--5
Unique Keyword 척도를 그대로 유지하되 — 네이버가 키워드 선택을 최대 5개로
제한하기 때문이다 — 이를 주 단위가 아닌 개별 리뷰 수준에서 무엇이 만들어내는지를
묻는다.

\subsect{2.2 리뷰 콘텐츠와 유용성}
직접적으로 관련된 또 다른 연구 흐름은, 리뷰 집계가 플랫폼 트래픽을 움직이는
요인이 아니라, 개별 리뷰 자체가 다른 독자에게 \emph{유용한가}를 묻는다. Wang,
Karmakar, Wang과 Pei(2025)는 약 160만 건의 아마존 리뷰를 활용하여, 리뷰
콘텐츠가 그 맥락(제품 설명 및 선행 리뷰)과 얼마나 다른가(비유사성)가 인지된
유용성에 미치는 영향을 연구하였다. 이들은 이 분야의 결과가 혼재되어 있다고
보고한다. Fresneda와 Gefen(2019; Wang 외, 2025에서 재인용)은 리뷰 내 고유
단어 수와 유용성 간의 \emph{긍정적} 관계를 발견한 반면, Namvar와 Chua(2023;
Wang 외, 2025에서 재인용)는 선행 리뷰와의 어휘적 비유사성이 유용성에
\emph{부정적} 영향을 미침을 발견하였다. 이러한 상반된 결과는, 본 보고서에서
리뷰가 받는 반응(아래 Model 4)을 다양성과 자동으로 함께 움직인다고 가정하지
않고 별도의 종속변수로 다루는 이유 중 하나이다.

\subsect{2.3 리뷰 감성과 독자 반응}
본 보고서에서 감성 점수는 BERT 계열 한국어 감성 분류기(hun3359/klue-bert-base-sentiment)로
산출하였다. BERT 기반 감성 분석 모형은 소비자 리뷰 맥락에서 기존 접근법보다
상당히 우수한 성능을 보이며, 제품 추천 맥락에서 91\%의 정확도를 기록한 바 있다
(Karabila, Darraz, EL-Ansari, Alami \& EL Mallahi, 2024). 다만 분류기 자체를
넘어, 본 보고서(IV장)와 이민영 외(2026)가 공통적으로 관측한 "더 긍정적인
리뷰일수록 독자 반응이 적다"는 패턴을 \emph{왜} 보이는지에 대한 이론적 근거는
본 연구를 위해 확인된 출처에서 다루지 않는다.

\todo{긍정적 감성이 독자 참여를 오히려 감소시키는 이유(예: 입소문에서의
부정성 편향, 부정적 정보의 진단적 가치)에 대한 이론적 근거가 추가로 필요하다.
검색어: "negativity bias online reviews helpfulness" 또는
"eWOM valence reader engagement".}

\subsect{2.4 연구 공백}
선행연구를 종합하면, 정보 다양성이 성과에 중요하다는 점(이민영 외, 2026)과
리뷰 콘텐츠가 인지된 유용성을 형성한다는 점(Wang 외, 2025)은 확인되었으나,
어느 쪽도 개별 리뷰 수준에서 다양성의 \emph{선행 요인} — 어떤 방문객이 어떤
상황에서 실제로 유익한 리뷰를 작성하는가 — 을 관측하지 않는다. 네이버는 모든
리뷰에 대해 동반자, 방문 목적, 대기 시간을 기록하고 있다. 본 보고서는 저자가
확인한 범위 내에서, 이러한 방문 맥락을 정보 다양성, 재방문 행동, 독자 반응과
단일한 리뷰 단위 틀 안에서 연결한 최초의 시도이다.

\sect{III. 연구방법}

\subsect{3.1 데이터}
한국 관광지에 대한 네이버 플레이스 리뷰를 관광지별로 개별 수집하였다. 데이터셋은
관광지 29곳, 작성자 8,276명, 리뷰 14,401건으로 구성되며, 이 중 11,165건은
텍스트를, 4,651건은 별점을 포함한다. 분석 단위는 업체·주가 아니라 관광지
$j$에 대한 개별 리뷰 $i$이다. 관광지는 사회적 평판(야놀자 리서치 "Korea's Top
500")과 내·외국인별 입장객 통계(한국관광공사)를 교차하여, 유형별 층(궁궐,
전통 지구, 시장, 유흥, 박물관, 현대 랜드마크, 자연) 내에서 선정하였다.

\subsect{3.2 변수}
네 개의 종속변수와, 네이버가 기록하는 방문 맥락 변수, 그리고 누가 언제
작성했는지에서 비롯되는 효과를 행동적 효과와 분리하기 위한 통제변수로
구성된다.

\tblVariables

\subsect{3.3 기술통계}
\tblDescriptives

키워드 수는 이민영 외(2026)가 보고한 0--5 척도와 동일한데, 이는 네이버가
키워드 선택을 최대 5개로 제한하기 때문이다. 별점은 거의 변동하지 않으며
(중앙값 5), 이것이 Model 3을 주요 결과가 아닌 탐색적 분석으로 보고하는
이유이다.

\subsect{3.4 상관분석}
\tblCorrelations

아래 회귀분석에 문제가 될 만큼 강한 상관은 없다. 가장 큰 값인 별점과 감성
점수의 상관(0.514)은 문제가 아니라, 감성 분류기가 실제 별점을 잘 따라가고
있음을 보여주는 근거이다. 키워드 수는 리뷰 경과일과 $-0.395$의 상관을 보이는데,
오래된 리뷰일수록 키워드가 적다는 것은 관측 기간 동안 키워드 사용이
증가했음을 의미하며, 이는 이민영 외(2026)가 사건으로 활용한 키워드 확장과
동일한 현상이다.

\subsect{3.5 연구모형}
관광지 고정효과 $\delta_j$를 포함한 리뷰 단위 회귀분석이며, 동일 작성자가
여러 리뷰를 작성하므로 표준오차는 작성자 단위로 군집화한다. 연결함수는 각
종속변수의 측정 척도에 따라 결정된다.

\vspace{-4pt}
\begin{align*}
  \text{M1 (재방문):}\quad &
    \operatorname{logit} P(\text{재방문}_{ij}) = \beta_0 + \mathbf{x}_{ij}^\top\boldsymbol\beta + \delta_j \\[2pt]
  \text{M2 (키워드 수):}\quad &
    P(\text{키워드 수}_{ij} \le k) = \big[1 + e^{-(\tau_k - \mathbf{x}_{ij}^\top\boldsymbol\beta - \delta_j)}\big]^{-1},\ k = 0,\dots,4 \\[2pt]
  \text{M3 (별점):}\quad &
    P(\text{별점}_{ij} \le k) = \big[1 + e^{-(\tau_k - \mathbf{x}_{ij}^\top\boldsymbol\beta - \delta_j)}\big]^{-1},\ k = 1,\dots,4 \\[2pt]
  \text{M4 (반응 수):}\quad &
    \text{반응 수}_{ij} \sim \text{NegBin}(\mu_{ij}),\quad
    \ln \mu_{ij} = \beta_0 + \mathbf{x}_{ij}^\top\boldsymbol\beta + \gamma\ln(\text{경과일}_{ij}) + \delta_j
\end{align*}
\vspace{-6pt}

M4는 포아송이 아닌 음이항 연결함수를 사용하는데, 본 표본에서 반응 수의
분산이 평균의 약 37배에 달하기 때문이다. 포아송 모형을 먼저 추정하여
이를 근거로 기각하였으며, 음이항 모형의 산포모수는 $\alpha = 2.44$로
추정되었다. 각 주 모형은 작성자별·관광지별 리뷰 1건만 무작위로 추출한
표본("작성자×관광지당 1건" 민감도 검정)에서 재추정하며, 전체 표본과
민감도 표본 양쪽에서 통과한 효과만 IV장에서 강건한 것으로 보고한다.
교차검증(AUC, RMSE)은 작성자 단위로 5-겹 분할하여, 동일 작성자가 학습과
검증에 함께 포함되지 않도록 하였다.

\tblVIF

\sect{IV. 연구결과}

\subsect{4.1 재방문 (Model 1)}
Model 1은 어떤 방문 맥락이 재방문 시 작성된 리뷰를 첫 방문 시 작성된
리뷰와 구별하는지를 묻는다. 교차검증 AUC는 0.773이며(계절·요일·인증
수단·팔로워 수를 추가한 확장 모형 Model 1x는 0.785; AIC는 9,596에서
9,372로 개선된다).

\begin{table}[H]\centering
\caption{Model 1: P(재방문), 로지스틱 회귀, 관광지 고정효과\\
\footnotesize $N = 11{,}165$건 (텍스트 포함, 작성자 6,070명), 재방문율 $= 22.6\%$, pseudo-$R^2 = 0.203$}
\footnotesize
\begin{tabular}{@{}lrrr@{}}
\toprule
변수 & OR & 95\% CI & $p$ \\
\midrule
대기 후 입장              & 0.51 & [0.35, 0.72] & 0.0002 *** \\
동반자: 혼자                & 1.68 & [1.18, 2.39] & 0.0037 ** \\
동반자: 연인/배우자           & 0.74 & [0.53, 1.04] & 0.0807 . \\
동반자: 친구                 & 0.58 & [0.42, 0.80] & 0.0011 ** \\
목적: 여행                    & 0.56 & [0.42, 0.75] & 0.0001 *** \\
목적: 데이트                   & 0.72 & [0.53, 0.97] & 0.0284 * \\
감성 점수 ($-1 \sim 1$)       & 0.70 & [0.60, 0.81] & 0.0000 *** \\
키워드 수 (0--5)               & 1.14 & [1.07, 1.23] & 0.0002 *** \\
ln(텍스트 길이)                 & 1.29 & [1.14, 1.47] & 0.0001 *** \\
ln(1 + 작성자 리뷰 수)           & 1.52 & [1.40, 1.65] & 0.0000 *** \\
\bottomrule
\end{tabular}
\\[3pt]
{\scriptsize 예약, 방문 정보 제공 및 나머지 동반자/목적 더미는 $p<.05$에서 유의하지 않았다. 전체 결과: \texttt{data/analysis/naver/model1\_revisit.csv}.}
\end{table}

\textbf{민감도 검정에서 강건한 효과}(작성자×관광지당 1건 표본에서 $p<.05$로
동일한 결론): 대기 후 입장, 혼자, 친구, 감성 점수, 키워드 수, ln(1 + 작성자
리뷰 수). \textbf{강건하지 않은 효과}: 방문 정보 제공, 연인/배우자, 목적:
데이트, ln(텍스트 길이) — 이들은 반복 작성자의 가중치를 낮추면 유의성을
잃으므로 해석에 주의가 필요하다.

확장 모형(Model 1x)에서 표 전체를 통틀어 가장 큰 효과는 \textbf{영수증
인증}이다(OR $= 0.35$, $p<.001$). 영수증으로 인증된 리뷰는 재방문자의
리뷰일 가능성이 현저히 낮다. ln(1 + 작성자 팔로워 수) 역시 긍정적이며(OR
$=1.18$, $p<.001$), 물놀이 관련 키워드는 재방문 확률과 부정적으로
연관된다(OR $=0.62$, $p=.009$).

\subsect{4.2 키워드 다양성 (Model 2)}
Model 2는 어떤 작성자와 방문 맥락이 더 다양한 리뷰를 생산하는지를 묻는다.
교차검증 RMSE는 1.366으로, 평균 예측 기준값 1.883보다 낮다(확장 모형
Model 2x는 1.337).

\begin{table}[H]\centering
\caption{Model 2: 키워드 수 (0--5), 순서형 로지스틱 회귀, 관광지 고정효과\\
\footnotesize $N = 14{,}401$, 평균 $= 2.97$개}
\footnotesize
\begin{tabular}{@{}lrrr@{}}
\toprule
변수 & OR & 95\% CI & $p$ \\
\midrule
텍스트 작성 (15자 이상)        & 4.02 & [3.51, 4.60] & 0.0000 *** \\
방문 정보 제공                  & 2.53 & [2.01, 3.18] & 0.0000 *** \\
재방문자 (2회 이상)              & 1.60 & [1.32, 1.93] & 0.0000 *** \\
동반자: 혼자                     & 1.67 & [1.27, 2.20] & 0.0003 *** \\
목적: 나들이                      & 1.85 & [1.60, 2.13] & 0.0000 *** \\
목적: 여행                         & 1.73 & [1.49, 2.02] & 0.0000 *** \\
목적: 친목                          & 1.59 & [1.31, 1.93] & 0.0000 *** \\
대기 후 입장                        & 0.75 & [0.64, 0.88] & 0.0004 *** \\
ln(1 + 작성자 리뷰 수)               & 1.17 & [1.13, 1.22] & 0.0000 *** \\
\bottomrule
\end{tabular}
\\[3pt]
{\scriptsize 모든 동반자·목적 더미는 방향성이 긍정적이며, 지면상 $p<.001$인 항목만 제시하였다. 전체 결과: \texttt{data/analysis/naver/model2\_keywords.csv}.}
\end{table}

\textbf{여러 설정에서 강건한 효과}(순서형 로지스틱과 이항 분포 강건성
검정 Model 2b 양쪽에서 유의): 재방문자, 텍스트 작성, 방문 정보 제공, 대기
후 입장, 혼자, 연인/배우자, 부모/가족, 목적: 여행·나들이·일상·데이트·친목,
ln(1 + 작성자 리뷰 수). 텍스트가 있는 리뷰만의 모형(Model 2t)에서는 감성
점수가 \emph{긍정적}으로 작용하고(OR $=2.26$, $p<.001$) ln(텍스트 길이)도
강하게 작용한다(OR $=1.86$, $p<.001$) — 더 길고 긍정적인 텍스트에 자연히
더 많은 키워드가 동반되기 때문이며, 바로 이 때문에 감성과 길이를 주 모형
(Model 2)에서는 제외하였다. 이들은 키워드 태그와 동일한 작성 행위의
산물이지, 그 선행 요인이 아니다.

\subsect{4.3 별점 (Model 3): 탐색적 분석}
본 코퍼스에서 별점은 거의 변동하지 않으며(중앙값 5, 5점 비율 73.9\%),
이에 따라 Model 3은 탐색적 분석으로 보고한다. McFadden $R^2 = 0.095$
($N = 4{,}651$). 기본, 이항(Model 3b), 텍스트 전용(Model 3t) 모형 전반에서
크고 일관된 유일한 효과는 \textbf{키워드 수}이다(기본 모형 OR $= 1.56$,
$p<.001$). 키워드를 더 많이 태그한 작성자는 별점도 더 높게 준다.
텍스트 전용 모형에서는 감성 점수가 다른 모든 변수를 압도하는데(OR $=
4.90$, $p<.001$), 이는 별점과 동일한 텍스트에서 감성을 산출했기 때문에
예상된 결과이며, 행동적 발견이라기보다 감성 분류기에 대한 검증으로
해석해야 한다. 텍스트 작성은 별점과 부정적으로 연관된다(OR $=0.64$,
$p<.001$) — 텍스트를 작성한 작성자는 키워드만 남긴 작성자보다 별점이
다소 낮다.

\subsect{4.4 반응 수 (Model 4)}
Model 4는 어떤 리뷰가 다른 이용자의 반응을 얻는지를 묻는다. 반응 발생
여부를 예측하는 교차검증 AUC는 계수형 모형의 예측 평균 기준 0.815이며
(전용 이항 모형 Model 4b는 0.827), pseudo-$R^2 = 0.212$이다.

\begin{table}[H]\centering
\caption{Model 4: 리뷰당 반응 수, 음이항 회귀, 관광지 고정효과\\
\footnotesize $N = 11{,}165$건 (텍스트 포함), 반응 1건 이상 $= 23.5\%$, 평균 $= 1.82$, $\alpha = 2.44$, pseudo-$R^2 = 0.212$}
\footnotesize
\begin{tabular}{@{}lrrr@{}}
\toprule
변수 & IRR & 95\% CI & $p$ \\
\midrule
ln(1 + 작성자 팔로워 수)      & 1.97 & [1.79, 2.18] & 0.0000 *** \\
ln(텍스트 길이)                & 1.64 & [1.44, 1.87] & 0.0000 *** \\
키워드 수 (0--5)                & 1.12 & [1.05, 1.20] & 0.0011 ** \\
\textbf{감성 점수 ($-1 \sim 1$)} & \textbf{0.77} & [0.64, 0.91] & 0.0032 ** \\
방문 정보 제공                   & 1.70 & [1.08, 2.68] & 0.0223 * \\
목적: 나들이                       & 0.63 & [0.45, 0.88] & 0.0070 ** \\
목적: 여행                          & 0.62 & [0.46, 0.83] & 0.0016 ** \\
방문: 가을                            & 0.71 & [0.61, 0.83] & 0.0000 *** \\
ln(1 + 작성자 리뷰 수)                & 0.90 & [0.81, 0.99] & 0.0378 * \\
\bottomrule
\end{tabular}
\\[3pt]
{\scriptsize 동반자 더미는 $p<.05$에서 유의하지 않았다. 전체 결과: \texttt{data/analysis/naver/model4\_reactions.csv}.}
\end{table}

\textbf{이는 이민영 외(2026)를 재현한다.} 해당 연구 역시 네이버의 키워드
확장 이전에는 감성이 참여(클릭 수)에 부정적으로 작용함을 보고하였다
($\alpha_{\text{Sentiment}} = -0.011 \sim -0.018$). 개별 리뷰 수준에서도
더 긍정적인 리뷰일수록 반응이 측정 가능할 만큼 적다 — 과도한 칭찬은 독자의
반응을 이끌어내지 못한다. \textbf{여러 설정에서 강건한 효과}(음이항 모형과
이항 반응-유무 모형 Model 4b 양쪽에서 유의): 키워드 수, ln(텍스트 길이),
방문: 가을, ln(1 + 작성자 팔로워 수). 감성 점수와 ln(1 + 작성자 리뷰 수)는
Model 4에서는 유의하나 Model 4b에서는 유의하지 않으므로, 확정적이기보다
시사적인 결과로 해석해야 한다.

\subsect{4.5 모형 성능 요약}
\begin{table}[H]\centering
\caption{교차검증 성능 (작성자 단위 5-겹 분할)}
\footnotesize
\begin{tabular}{@{}llll@{}}
\toprule
모형 & 지표 & 값 & 기준 \\
\midrule
M1 재방문                 & AUC          & 0.773 (확장 0.785) & -- \\
M2 키워드 수                & RMSE          & 1.366 (확장 1.337)   & 1.883 (평균 예측) \\
M3 별점                     & --            & 탐색적                & -- \\
M4 반응 수                   & AUC (반응 유무) & 0.815                 & -- \\
M4b 반응 유무                 & AUC            & 0.827                 & -- \\
\bottomrule
\end{tabular}
\end{table}

\sect{V. 결론 및 시사점}

\subsect{5.1 연구결과 요약}
\begin{itemize}[leftmargin=14pt, itemsep=2pt]
  \item 네이버 플레이스의 정보 다양성(키워드 수)은 방문 유형보다, 작성자가
        텍스트를 썼는지, 방문 정보를 제공했는지, 이전에 리뷰를 자주
        작성했는지 등 플랫폼 참여도를 따른다.
  \item 감성 점수는 독자 참여(반응)에 부정적으로 작용하며, 이는 이민영
        외(2026)를 업체·주 단위가 아닌 개별 리뷰 수준에서 재현한 결과이다.
  \item 대기(입장 전 기다림)는 재방문에 대한 가장 강한 부정적 신호이며,
        본 연구 결과 중 관광지 운영자가 직접 개선할 수 있는 유일한 요인이다.
  \item 본 코퍼스에서 별점은 상대적으로 정보량이 적으며, 키워드 태그가
        더 정보력 있는 종속변수이다.
\end{itemize}

\subsect{5.2 시사점}
플랫폼 측면에서 정보 다양성은 방문 유형보다 작성자의 참여도를 따르므로,
다양한 키워드를 유도하는 기능은 처음 방문한 이용자보다 경험이 많은
작성자에게 더 효과적일 것이다. 관광지 운영 측면에서는, 본 연구 결과 중
대기 시간 단축이 방문객의 재방문 여부에 실질적으로 영향을 미칠 수 있는
유일한 수단이다.

\subsect{5.3 한계 및 향후 과제}
\begin{itemize}[leftmargin=14pt, itemsep=2pt]
  \item 결과는 상관관계이다. 이민영 외(2026)는 네이버의 키워드 확장 사건을
        활용해 효과를 식별하고 모든 독립변수를 $t-1$로 처리하지만, 본
        설계는 횡단면이며 두 장치를 모두 갖추지 못했다.
  \item 관광지 고정효과만 포함하였으나, 이민영 외(2026)는 업체와 시간의
        양방향 고정효과를 사용한다. 월 고정효과 추가와, 고정효과 모형의
        타당성을 검정하는 공식적인 Hausman 검정이 다음 과제이다.
  \item Model 1은 리뷰 작성 이전에 이미 발생한 방문을 사후에 작성된
        리뷰로 예측하므로, 재방문의 원인을 설명하기보다 재방문자의 특성을
        식별하는 모형이다.
  \item 단일 플랫폼 자료이므로 네이버 작성자 집단과 한국인 작성자 전반을
        구분할 수 없다.
\end{itemize}

\vspace{8pt}
\subsect{참고문헌}
\begin{itemize}[leftmargin=14pt, itemsep=3pt]
  \item Bae, E., Chang, M., Park, E., \& Kim, D. (2017). The effect of
        Hallyu on tourism in Korea. \textit{Journal of Open Innovation:
        Technology, Market, and Complexity}, 3, 22.
  \item Chang, S. (2021). Food tourism in Korea. \textit{Journal of
        Vacation Marketing}.
  \item 이민영·정동영·김지영 (2026). 정보 다양성이 온라인 리뷰와 업체
        성과에 미치는 영향. \textit{소비자학연구}, 37(1).
        \url{https://doi.org/10.35736/JCS.37.1.1}
  \item Karabila, I., Darraz, N., EL-Ansari, A., Alami, N., \& EL Mallahi,
        M. (2024). BERT-enhanced sentiment analysis for personalized
        e-commerce recommendations. \textit{Multimedia Tools and
        Applications}, 83, 56463--56488.
  \item Ko, S., Kang, S., Kang, H., \& Lee, M. J. (2018). An exploration of
        foreign tourists' perceptions of Korean food tour: a
        factor-cluster segmentation approach. \textit{Asia Pacific Journal
        of Tourism Research}, 23(8), 833--846.
  \item Wang, S., Karmakar, S., Wang, F., \& Pei, Y. (2025). Content
        dissimilarity and online review helpfulness: Contextual insights.
        \textit{Journal of Business Research}, 187, 115068.
\end{itemize}

\vspace{4pt}
{\footnotesize\textit{위 참고문헌 중 Fresneda \& Gefen(2019), Namvar \&
Chua(2023), Dellarocas(2003)는 각각 Wang 외(2025)와 이민영 외(2026) 안에서
재인용된 것으로, 본 연구를 위해 원문을 직접 확인하지 않았다. 이 보고서를
제출하기 전에 반드시 원문을 대조해야 한다. 1.1절, 2.3절, 2.4절의
[SOURCE NEEDED] 표시 세 곳은 추가로 필요한 문헌을 구체적으로 명시한
것이며, 가상의 인용을 대신한 자리표시자가 아니다.}}

\end{document}
